\scanpage{22}
\mspage{11}
un syst. transitif d'iso. entre les $\pi_t$ $(t\in\unsure{\tau|\mathcal T})$, permettant de les identifier tous à un $=T$. Ainsi $\langle\unsure{\tau|\mathcal T},\pi\rangle$ devient un $T$-groupoïde, c.àd. un groupoïde (pas néc. connexe) où les grpes fondamentaux des comp. conn. sont comm. et tous identifiés à $T$.
De plus, $\unsure{\mathcal E|\Sigma}$ opère sur ce $T$-groupoïde (en opérant également sur $T$) via ses opérations sur le système $(\unsure{\tau|\mathcal T},\pi)$ ($\pi$-espace à gauche) en opérant sur $\unsure{\tau|\mathcal T}$ par l'op. donnée et sur $\pi$ via autom int. de $\unsure{\mathcal E|\Sigma}$, enfin sur $T$ via $\mathcal G$ via $\chi\ldots.$

[Si $\unsure{\mathcal E|\Sigma}$ opère sur $\unsure{\tau|\mathcal T}$, restreignant cette op. à $\pi$ on trouve un groupoïde $\langle\unsure{\tau|\mathcal T},\pi\rangle$ sur lequel $\unsure{\mathcal E|\Sigma}$ opère. Supposant ce groupoïde abélien, on a donc une famille de grpes abéliens \inserted{$(T_j)_{j\in J}$} indexée par $J=\unsure{\tau|\mathcal T}/\pi$, sur laquelle opère $\unsure{\mathcal E|\Sigma}$ via $\unsure{\mathcal E|\Sigma}/\pi=\mathcal G$. Ceci dit, on suppose donné une trivialisation de cette famille $(T_j)$ par $T\simeq T_j$, compatible avec op. de
\mspage{21}
